% !TeX root = ../main.tex
\section{Potenze a esponente intero positivo}
% Fonte: Lezione 6, pp.~14--15

Sia $n\in\N^*$. La potenza $x^n$ è il prodotto di $n$ fattori uguali a $x$ ed è definita per ogni $x\in\R$. Il comportamento dipende dalla parità dell'esponente.
\begin{itemize}
  \item \textbf{Esponente pari.} Se $n$ è pari, $x^n$ è non negativa su $\R$, nulla soltanto in $0$ e strettamente positiva per $x\ne0$. È una funzione pari, strettamente decrescente su $(-\infty,0]$ e strettamente crescente su $[0,+\infty)$. La sua immagine è $[0,+\infty)$:
  \[
  x\longmapsto x^n\colon\R\longrightarrow[0,+\infty).
  \]
  \item \textbf{Esponente dispari.} Se $n$ è dispari, $x^n$ ha lo stesso segno di $x$: è non negativa su $[0,+\infty)$ e non positiva su $(-\infty,0]$, con segno stretto fuori da $0$. È dispari e strettamente crescente su tutto $\R$. La sua immagine è $\R$:
  \[
  x\longmapsto x^n\colon\R\longrightarrow\R.
  \]
\end{itemize}

\begin{proof}[Verifica del segno e della monotonia]
\textbf{Segno e simmetria.} Seguono dal numero di fattori negativi e dall'identità $(-x)^n=(-1)^n x^n$.

\textbf{Crescita sui non negativi.} Per $0\le u<v$, la fattorizzazione
\[
v^n-u^n=(v-u)\sum_{k=0}^{n-1}v^{n-1-k}u^k>0
\]
prova la stretta crescita sulla semiretta positiva; per $n=1$ la somma vale $1$.

\textbf{Semiretta negativa e confronto fra segni opposti.} Per $u<v\le0$, si ha $0\le-v<-u$. La crescita sui non negativi, insieme alla simmetria, dà $u^n>v^n$ nel caso pari e $u^n<v^n$ nel caso dispari. Nel caso dispari, punti di segno opposto hanno anche immagini di segno opposto: la crescita vale su tutta la retta.

Le immagini dichiarate seguono dall'esistenza delle radici, dimostrata nella prossima sezione.
\end{proof}

\section{Radici e potenze a esponente razionale}
% Fonte: Lezione 6, pp.~16--17

Per ogni $n\in\N^*$ la restrizione della potenza alla semiretta $[0,+\infty)$ è biettiva su $[0,+\infty)$. La sua inversa si chiama \emph{radice $n$-esima}:
\[
\sqrt[n]{\,\cdot\,}\colon[0,+\infty)\longrightarrow[0,+\infty),
\qquad \sqrt[n]{x}=x^{1/n}.
\]
Più precisamente, $\sqrt[n]{x}$ è l'unico numero \emph{non negativo} la cui potenza $n$-esima è $x$; è strettamente positivo quando $x>0$, mentre $\sqrt[n]{0}=0$.

\begin{proposizione}[Esistenza e unicità della radice]
Per $a\ge0$ e $n\in\N^*$ esiste uno e un solo $c\ge0$ tale che $c^n=a$.
\end{proposizione}
\begin{proof}
\textbf{Esistenza: scelta del candidato.} Per $a=0$ basta $c=0$. Per $a>0$ consideriamo
\[
S=\{t\ge0:t^n\le a\}.
\]
È non vuoto perché contiene $0$ ed è limitato superiormente: se $t>\max\{1,a\}$, allora $t^n\ge t>a$. Esiste dunque $c=\sup S\ge0$.

\textbf{Stima della variazione.} Poniamo $B=c+1$ e $K=nB^{n-1}>0$. La fattorizzazione della differenza di potenze dà, per $u,v\in[0,B]$,
\[
|u^n-v^n|\le K|u-v|,
\]
perché ciascuno dei $n$ prodotti della somma è al più $B^{n-1}$.

\textbf{Caso}\enspace\passomath{c^n<a} Scegliamo $0<h<\min\{1,(a-c^n)/K\}$. Allora $(c+h)^n\le c^n+Kh<a$, perciò $c+h\in S$, in contraddizione con il fatto che $c$ è un maggiorante.

\textbf{Caso}\enspace\passomath{c^n>a} In questo caso $c>0$; scegliamo $0<h<\min\{c,(c^n-a)/K\}$. Si ha $(c-h)^n\ge c^n-Kh>a$. Per la stretta crescita della potenza sui non negativi, nessun $t\ge c-h$ appartiene a $S$. Quindi $c-h$ sarebbe un maggiorante di $S$ più piccolo di $c$, ancora impossibile. Resta necessariamente $c^n=a$.

\textbf{Unicità.} La potenza è strettamente crescente sui non negativi: due numeri non negativi distinti non possono avere la stessa potenza $n$-esima. Questo vale anche per $a=0$.
\end{proof}

\textbf{Indice dispari.} Se $n$ è dispari, la potenza è invertibile su tutta la retta e la radice si estende a una funzione $\R\to\R$ ponendo, per $x<0$,
\[
\sqrt[n]{x}=-\sqrt[n]{-x}.
\]
È strettamente crescente e dispari.

\textbf{Indice pari.} Se $n$ è pari, nessun reale ha potenza $n$-esima negativa, quindi il dominio reale della radice è soltanto $[0,+\infty)$. Per $a>0$, l'equazione $t^n=a$ con $n$ pari ha due soluzioni, $\pm\sqrt[n]{a}$: il simbolo di radice sceglie quella non negativa. In particolare
\[
(\sqrt[n]{x})^n=x\quad(x\ge0),\qquad
\sqrt[n]{x^n}=\begin{cases}|x|&n\text{ pari},\\x&n\text{ dispari}.
\end{cases}
\]

\begin{schema}
\begin{tikzpicture}[x=1.35cm,y=.9cm]
  \draw[asse] (-1.7,0) -- (1.8,0) node[right] {$x$};
  \draw[asse] (0,-2.8) -- (0,3) node[above] {$y$};
  \draw[intervallo,domain=-1.55:1.55,samples=90] plot (\x,{\x*\x});
  \draw[thick,orange!80!black,domain=-1.38:1.38,samples=90]
    plot (\x,{\x*\x*\x});
  \node[coloreschema,left,fill=white,inner sep=1pt] at (-1.2,1.9) {$x^2$};
  \node[orange!80!black,right] at (1.1,1.2) {$x^3$};
  \node[below] at (0,-3) {Potenze: esponente pari e dispari};
\end{tikzpicture}
\hspace{.9cm}
\begin{tikzpicture}[x=1.15cm,y=1.25cm]
  \draw[asse] (-2.3,0) -- (2.5,0) node[right] {$x$};
  \draw[asse] (0,-1.8) -- (0,2.15) node[above] {$y$};
  \draw[intervallo,domain=0:1.5,samples=75] plot ({\x*\x},\x);
  \draw[thick,orange!80!black,domain=-1.3:1.3,samples=90]
    plot ({\x*\x*\x},\x);
  \node[coloreschema,above left] at (1.9,1.45) {$\sqrt x$};
  \node[orange!80!black,below right] at (1.4,1.1) {$\sqrt[3]x$};
  \node[below] at (0,-2.16) {Radici: indice pari e dispari};
\end{tikzpicture}
\end{schema}

\textbf{Potenze razionali positive.} Combinando potenze intere e radici si definiscono le potenze a esponente razionale positivo. Per $p,q\in\N^*$ e $x\ge0$ poniamo
\[
x^{p/q}=(\sqrt[q]{x})^p=\sqrt[q]{x^p}.
\]
\textbf{Compatibilità delle due scritture.} Le due espressioni coincidono perché sono non negative e, elevate alla $q$-esima potenza, danno entrambe $x^p$.

\textbf{Indipendenza dalla frazione.} Se due frazioni rappresentano lo stesso razionale, cioè se $p/q=p'/q'$, le due potenze ottenute, elevate alla potenza $qq'$, danno rispettivamente $x^{pq'}$ e $x^{p'q}$, uguali poiché $pq'=p'q$. La stretta crescita della potenza $qq'$ sui non negativi conclude.

Per basi negative occorre verificare separatamente il dominio reale delle radici coinvolte; la definizione appena data riguarda esplicitamente le basi non negative.
